\documentclass[11pt]{article}
\usepackage[margin=1in]{geometry}
\usepackage[T1]{fontenc}
\usepackage{lmodern}
\usepackage{microtype}
\usepackage{amsmath,amssymb}
\usepackage[dvipsnames]{xcolor}
\usepackage{enumitem}
\usepackage[colorlinks=true,linkcolor=MidnightBlue,urlcolor=MidnightBlue,citecolor=MidnightBlue]{hyperref}
\setlist[itemize]{leftmargin=1.4em,itemsep=0.3em,topsep=0.35em}
\setlength{\parindent}{0pt}
\setlength{\parskip}{0.62em}

\newcommand{\statusbox}[1]{%
  \begin{center}
  \fcolorbox{MidnightBlue}{MidnightBlue!6}{%
    \parbox{0.91\linewidth}{\centering\bfseries\color{MidnightBlue}#1}}
  \end{center}}

\title{Literature Status Note\\[0.25em]
\large Problems A--C from arXiv:1012.4514}
\author{Run \texttt{fa\_banach\_001} --- lane 5}
\date{27 August 2026}

\begin{document}
\maketitle

\statusbox{LITERATURE ALREADY ANSWERED --- PROBLEMS A, B, AND C}

\section*{Source questions}

Levy and Shalit's survey \emph{Dilation theory in finite dimensions} asks four
labeled questions \cite{LS}.  This note identifies exact later answers to the
first three.

\begin{itemize}
\item \textbf{Problem A} (source PDF page 7): can a commuting tuple of
contractive $3\times3$ matrices violate
\[
 \|p(T_1,\ldots,T_d)\|\leq \sup_{z\in\mathbb T^d}|p(z)|?
\]
\item \textbf{Problem B} (source PDF page 8): are Brehmer's positivity
inequalities equivalent to the existence, for every $N$, of a regular unitary
$N$-dilation on a finite-dimensional space?
\item \textbf{Problem C} (source PDF page 9): does every pair of commuting
contractions on a finite-dimensional space have a finite-dimensional unitary
$N$-dilation for each fixed $N$?
\end{itemize}

\section*{Exact literature answers}

\textbf{Problem A: no counterexample exists.}
Knese's Theorem 1.2, on supporting PDF page 2, states that the von Neumann
inequality holds for every $d$-tuple of commuting contractive $3\times3$
matrices \cite{Knese}.  This is precisely the universal negation of the
existence question in Problem A.  Knese attributes the decisive three-point
Pick interpolation input to Kosi\'nski and supplies the operator-theoretic
reductions.

\textbf{Problem B: yes.}
McCarthy and Shalit's Theorem 1.7 and Corollary 1.8, on supporting PDF page 4,
prove that a finite-dimensional commuting tuple has finite-dimensional regular
unitary $N$-dilations for every $N$ exactly when it has a regular unitary
dilation; Brehmer's inequalities characterize the latter \cite{MS}.  Their
corollary is word-for-word the equivalence requested in Problem B.

\textbf{Problem C: yes.}
McCarthy and Shalit's Theorem 1.2, stated on supporting PDF page 2, proves more
generally that whenever a commuting tuple on a finite-dimensional space has a
unitary dilation, it has a finite-dimensional unitary $N$-dilation for every
fixed $N$ \cite{MS}.  Apply Ando's dilation theorem to a pair of commuting
contractions.

\section*{Proof intuition for the finite-dimensional reduction}

The common mechanism behind Problems B and C is finite moment compression.
A commuting unitary dilation yields an operator-valued measure on the torus.
Only finitely many moments are needed at degree $N$.  A matrix-valued
Carath\'eodory argument replaces that measure by a finite atomic positive
operator-valued measure with the same moments, and finite-dimensional Naimark
dilation turns its atoms into commuting unitaries.  The regular version retains
both positive and negative Laurent moments.  This is the exact finite-dimensional
bridge that the source paper sought.

\section*{Scope boundary}

Problem D asks for a ``commutant lifting theorem'' for unitary $N$-dilations,
without fixing a unique formulation.  McCarthy and Shalit explicitly say that
the appropriate analogue is unclear and show that one natural fixed-dilation
version fails already for a particular Halmos $1$-dilation \cite[Sec.~2.3]{MS}.
That observation does not settle every plausible existential version.  We
therefore record Problems A--C as answered and make no durable status claim for
Problem D.

\section*{Review recommendation}

Classify the queued Problem-A signal as answered by Knese.  Record Problems B
and C as additional exact answers from McCarthy--Shalit.  Do not spend a new
proof attempt on this paper unless Problem D is first reformulated precisely.

\begin{thebibliography}{9}
\bibitem{LS}
E.~Levy and O.~M.~Shalit,
\emph{Dilation theory in finite dimensions: the possible, the impossible and
the unknown}, arXiv:1012.4514; Rocky Mountain J. Math. 44 (2014), 203--221.

\bibitem{Knese}
G.~Knese,
\emph{The von Neumann inequality for $3\times3$ matrices},
arXiv:1508.06488; Bull. Lond. Math. Soc. 48 (2016), 53--57.

\bibitem{MS}
J.~E.~McCarthy and O.~M.~Shalit,
\emph{Unitary $N$-dilations for tuples of commuting matrices},
arXiv:1105.2020; Proc. Amer. Math. Soc. 141 (2013), 563--571.
\end{thebibliography}

\end{document}
